\section{Model formulations}
\label{sec:models}

\subsection{Thermal balance with a distributed input delay}
The measured input $u(t)$ is heater power in percentage points, and the measured
output $T(t)$ is the first temperature sensor in degrees Celsius. We fixed the
heater conversion $c=\alpha P$ at $0.032\,\mathrm{W}$ per percentage point, using
the inherited values $\alpha=0.00016$ and $P=200$. The ambient temperature
$\Tamb$ was set to the first measured temperature in each experiment. This
choice conditions the analysis on an uncertain ambient estimate.

For a chain of $n$ delay states, the auxiliary variables $z_i(t)$ have the same
units as $u(t)$. The physical parameters $\Ua>0$, $\Cp>0$, and $\theta>0$
denote thermal conductance, effective heat capacity, and mean input-delay time,
respectively. The model is
\begin{align}
 \dot z_1 &= \frac{n}{\theta}(u-z_1), &
 \dot z_i &= \frac{n}{\theta}(z_{i-1}-z_i),\quad i=2,\ldots,n,\label{eq:chain}\\
 \Cp\dot T &= -\Ua(T-\Tamb)+c z_n.\label{eq:physical}
\end{align}
The input--output coordinates $K=c/\Ua$ and $\tau=\Cp/\Ua$ are steady-state
gain and thermal time constant. In these coordinates the temperature equation is
\begin{equation}
 \dot T=-\frac{T-\Tamb}{\tau}+\frac{K}{\tau}z_n.
 \label{eq:io}
\end{equation}
Both parameterizations describe the same positive model family. Their inverse
map is $\Ua=c/K$ and $\Cp=c\tau/K$. Distinct rectangular parameter bounds need
not map onto each other; our numerical comparison records active bounds and
compares common interior solutions.

For zero initial deviations, the Laplace variable $s$ gives the transfer function
\begin{equation}
 G_n(s)=\frac{K}{(1+\tau s)(1+\theta s/n)^n}.
 \label{eq:transfer}
\end{equation}
The delay chain has an Erlang impulse-response kernel with mean $\theta$ and
variance $\theta^2/n$; Appendix~\ref{sec:structural} derives these quantities.
The linear chain construction is consistent with the Erlang dwell-time
representation studied by \citet{hurtado-2019-generalizations}. A finite chain
responds before time $\theta$ and therefore represents a distributed lag. As
$n$ increases at fixed $\theta$, its transfer function tends to
$\exp(-\theta s)$. We also fitted the limiting pure-delay model,
\begin{equation}
 \tau\dot T(t)=-(T(t)-\Tamb)+K u(t-\theta).
 \label{eq:puredelay}
\end{equation}
Increasing $n$ changes the transfer-function shape and adds dynamic states;
it does not add estimated parameters to the three-parameter family.

\subsection{Initial-state assumptions}
We conditioned $T(0)$ on the first observation. The cold-start assumption set
all $z_i(0)=0$, while the input-equilibrated assumption set all $z_i(0)=u(0)$.
The pure-delay model used the corresponding zero or constant input prehistory.
These assumptions encode different unobserved histories of heater operation.
Neither was established by the recorded input before the first sample, because
that history is absent from the supplied data.

\subsection{Original two-state thermal model}
The inherited model describes heater temperature $T_{\mathrm H}$ and sensor
temperature $T_{\mathrm S}$. Its conductances $\Ua,\Ub>0$ and heat capacities
$\CpH,\CpS>0$ satisfy
\begin{align}
 \CpH\dot T_{\mathrm H}&=-\Ua(T_{\mathrm H}-\Tamb)
                       -\Ub(T_{\mathrm H}-T_{\mathrm S})+cu,\\
 \CpS\dot T_{\mathrm S}&=\Ub(T_{\mathrm H}-T_{\mathrm S}).
 \label{eq:original}
\end{align}
Only $T_{\mathrm S}$ was used in this single-output comparison. Both initial
temperatures were fixed at the first sensor reading. Under these conditions its
transfer function depends on three combinations of four physical parameters:
\begin{equation}
 G_{\mathrm{orig}}(s)=\frac{K}{1+a s+b s^2},\qquad
 K=\frac{c}{\Ua},\quad
 a=\frac{\CpH}{\Ua}+\frac{\CpS}{\Ub}+\frac{\CpS}{\Ua},\quad
 b=\frac{\CpH\CpS}{\Ua\Ub}.
 \label{eq:original-transfer}
\end{equation}
Appendix~\ref{sec:structural} exhibits the resulting continuum of equivalent
physical parameters. The one-stage cold-start model has $a=\tau+\theta$ and
$b=\tau\theta$, so its zero-state transfer family coincides with this original
model. The pole labels $\tau$ and $\theta$ can be interchanged. We used the
convention $\tau\geq\theta$ when the exchanged point remained within the
specified parameter bounds and retained both modes in likelihood profiles.
